\documentclass[10pt,oneside]{book}

% Page geometry — 125 × 200 mm (Gyldendal 1999 Tractatus format).
% Wide LEFT margin holds the proposition numbers (marginal numbering style):
% numbers are \llap'd into this margin so the body text block is a clean
% rectangle with all lines (including continuation) flush left at the body
% column. Right margin is narrow (no marginalia on the right).
% Top includes the running header band; footskip minimal (no footer text).
\usepackage[
  paperwidth=125mm,
  paperheight=200mm,
  top=20mm,
  bottom=18mm,
  left=28mm,
  right=12mm,
  headheight=12pt,
  headsep=6mm,
  footskip=8mm,
]{geometry}

% UTF-8 + Norwegian (nynorsk)
\usepackage{fontspec}
\usepackage{polyglossia}
\setdefaultlanguage[variant=nynorsk]{norwegian}

% Body font — the digibok is set in Sabon. EB Garamond is the closest
% free serif with the same general feel; fall back to TeX Gyre Pagella.
\IfFontExistsTF{EB Garamond}{
  \setmainfont{EB Garamond}[
    Numbers={Lining,Proportional},
    SmallCapsFeatures={LetterSpace=4},
  ]
}{
  \setmainfont{TeX Gyre Pagella}
}
% Body size 10.5 pt on 12.5 pt leading (close to Sabon 10.5/12 in the digibok)
\renewcommand{\normalsize}{\fontsize{10.5}{12.5}\selectfont}
\normalsize

% Better line breaking (must be loaded after font setup)
\usepackage{microtype}
\usepackage{endnotes}
\let\footnote=\endnote

% Math-symbol fallback font — EB Garamond lacks ∀ ∃ ∈ ⊆ ⊃ ≠ etc.,
% so the Python preprocessor wraps them in \msym{...} which uses
% Cambria Math (alltid present on Windows).
\newfontfamily{\mathfbfont}{Cambria Math}[Scale=0.95]
\newcommand{\msym}[1]{{\mathfbfont #1}}

% Headings — small caps, centred, generous spacing
\usepackage{titlesec}
\titleformat{\chapter}[block]
  {\normalfont\centering\scshape\addfontfeature{LetterSpace=10}\large}
  {}{0em}{\MakeLowercase}
\titlespacing*{\chapter}{0pt}{2.5em}{2.5em}

\titleformat{\section}[block]
  {\normalfont\centering\scshape\addfontfeature{LetterSpace=6}\normalsize}
  {}{0em}{\MakeLowercase}
\titlespacing*{\section}{0pt}{2.5em}{1.5em}

\titleformat{\subsection}[block]
  {\normalfont\scshape\addfontfeature{LetterSpace=2}\small}
  {}{0em}{\MakeLowercase}
\titlespacing*{\subsection}{0pt}{1.5em}{0.8em}

% Hide chapter numbers (chapters are titled by name)
\renewcommand{\thechapter}{}

% Running headers — chapter title at left, page number at right, thin rule.
% The header rule is extended LEFT by 16 mm via \fancyheadoffset so it
% spans the full content width INCLUDING the marginal number column
% (otherwise the rule looks narrow because it's anchored to the body
% text block which sits to the right of the number margin).
\usepackage{fancyhdr}
\pagestyle{fancy}
\fancyhf{}
\fancyhead[L]{\small\scshape\leftmark}  % chapter title at left on every page
\fancyhead[R]{\small\thepage}           % page number at right on every page
\fancyheadoffset[L]{16mm}               % extend header left over the number margin
\renewcommand{\headrulewidth}{0.4pt}
\renewcommand{\footrulewidth}{0pt}
% plain style (chapter-opening pages): identical to fancy (full header).
\fancypagestyle{plain}{
  \fancyhf{}
  \fancyhead[L]{\small\scshape\leftmark}
  \fancyhead[R]{\small\thepage}
  \fancyheadoffset[L]{16mm}
  \renewcommand{\headrulewidth}{0.4pt}
  \renewcommand{\footrulewidth}{0pt}
}
% title style — used ONLY on the title page: nothing in the header at all.
\fancypagestyle{titlepage}{
  \fancyhf{}
  \fancyheadoffset[L]{0pt}
  \renewcommand{\headrulewidth}{0pt}
  \renewcommand{\footrulewidth}{0pt}
}

% \silentchapter{display name}{lowercase mark}
%   Starts a new page, sets the running mark, adds a TOC entry, and uses the
%   plain pagestyle for the opening page — but DOES NOT emit a centred chapter
%   title (the chapter is identified by the running header and by its content).
\newcommand{\silentchapter}[2]{%
    \markboth{#2}{}%
  \addcontentsline{toc}{chapter}{#1}%
}

% Redefine \tableofcontents to skip the centred "Innhald" title and to
% emit the entire TOC at \footnotesize so it fits compactly on one page.
\makeatletter
\renewcommand\tableofcontents{%
  \markboth{innhald}{}%
  \begingroup\footnotesize
  \@starttoc{toc}%
  \endgroup
}
% Very compact chapter-level TOC entries
\renewcommand*\l@chapter[2]{%
  \ifnum \c@tocdepth >\m@ne
    \vskip 0pt%
    \setlength\@tempdima{1.5em}%
    \begingroup
      \footnotesize
      \parindent \z@ \rightskip \@pnumwidth
      \parfillskip -\@pnumwidth
      \leavevmode \normalfont
      \advance\leftskip\@tempdima
      \hskip -\leftskip
      #1\nobreak\leaders\hbox{$\m@th\mkern \@dotsep mu\hbox{.}\mkern \@dotsep mu$}\hfill
      \nobreak\hb@xt@\@pnumwidth{\hss #2}\par
    \endgroup
  \fi}
% Compact section / subsection entries — tight indents.
% Font size is set via \footnotesize on the entire TOC (see \tableofcontents
% redefinition above), so the entries inherit it.
\renewcommand*\l@section{\@dottedtocline{1}{1.0em}{2.0em}}
\renewcommand*\l@subsection{\@dottedtocline{2}{3.0em}{2.6em}}
\makeatother

% ─── Proposition environment ──────────────────────────────────────────────
%
% \prop{1.21}{d}{Body text..\footnote{Udelelegheit tyder her at vidare dekomponering øydelegg objektets kausale funksjon i den gjevne lyskjegla.}.}
%   ⇒  number sits in the LEFT MARGIN (\llap'd outside the text block);
%      body text fills the entire text block as a clean justified rectangle;
%      continuation lines wrap to the body column (= text margin).
%
% This is the Gyldendal 1999 Tractatus / Töpfer marginal-numbering layout:
% the body text is a single uniform column and the proposition numbers hang
% out into the white space to the left.
%
% Numbers in the digibok appear in the same weight as the body (regular,
% not bold). The status letter is rendered as an italic superscript right
% after the number.
\usepackage{calc}
\newlength{\propnumwidth}
\setlength{\propnumwidth}{13mm}  % column reserved for the number itself
\newlength{\propnumgap}
\setlength{\propnumgap}{2mm}     % gap between number column and body text

% \prop{num}{status}{body}
%   The number is placed in a left-aligned 13 mm box, then a 2 mm gap,
%   and the whole 15 mm slug is \llap'd into the LEFT margin so its right
%   edge sits flush against the text-block left edge. Body text starts at
%   the text margin and wraps there too — a perfect rectangle.
%   The number itself is rendered in a heavier weight so it stands out
%   against the regular body text.
\newcommand{\prop}[3]{%
  \par\addvspace{6pt}%
  \noindent\llap{\makebox[\propnumwidth][l]{\textbf{#1}\textsuperscript{\textit{#2}}}\hspace{\propnumgap}}%
  #3\par%
  \addvspace{2pt}%
}

% Continuation paragraph — plain paragraph at the body column (no number).
\newcommand{\propcont}[1]{%
  \par\addvspace{3pt}%
  \noindent #1\par%
}

% Italic transition between chapters
\newenvironment{overgang}{%
  \par\addvspace{12pt}%
  \begingroup\itshape\noindent
}{%
  \endgroup\par\addvspace{12pt}%
}

% Glossary entries — bold term, regular body, with breathing room
% between entries. Spans up to two pages comfortably.
\newenvironment{ordliste}{%
  \begingroup\small\setlength{\parskip}{2pt plus 0.6pt}%
  \linespread{0.97}\selectfont
}{%
  \endgroup
}
\newcommand{\ordlisteentry}[2]{%
  \par\addvspace{2.5pt}%
  \noindent\textbf{#1:} #2\par%
}

% Compact Føreord body — small font, very tight parskip and linespread
% so the foreword fits on a single page.
\newenvironment{foreordbody}{%
  \begingroup\small\setlength{\parskip}{0.5pt plus 0.2pt}%
  \linespread{0.90}\selectfont
}{%
  \endgroup
}

% A.6.x explanatory body text — slightly smaller, tight line spacing
% so each A.6.x entry (heading + figure + explanation) fits compactly.
\newcommand{\anote}[1]{%
  \par\addvspace{1pt}%
  {\footnotesize\linespread{0.93}\selectfont\noindent #1\par}%
  \addvspace{1pt}%
}

% Reference entries — very compact, tight line spacing, hanging indent.
% Designed to fit ~30 entries on a single page.
\newcommand{\refentry}[1]{%
  \par\addvspace{0pt}%
  \begingroup\footnotesize\linespread{0.92}\selectfont
  \noindent\hangindent=4mm\hangafter=1 #1\par%
  \endgroup
}

% Figures — centred, max width = text width
\usepackage{graphicx}
\graphicspath{{analysis/figures/}}
\usepackage{caption}
\captionsetup{
  font={small,it},
  labelformat=empty,
  justification=centering,
  skip=4pt,
}
\usepackage{float}

% Formula display blocks (for D4, D6, etc.) — monospace, small left indent
\newenvironment{formelblokk}{%
  \par\addvspace{4pt}%
  \begingroup\ttfamily\small\leftskip=4mm\noindent
}{%
  \endgroup\par\addvspace{4pt}%
}

% Avoid widows/orphans
\widowpenalty=10000
\clubpenalty=10000

% Body paragraph defaults — tight for the small page
\setlength{\parindent}{0pt}
\setlength{\parskip}{3pt plus 0.6pt minus 0.4pt}
\linespread{1.0}

% Wide environment for back-matter sections (Etterord, Referansar) that
% should fill the full content width including the 16 mm marginal number
% column. Pulls the left edge 16 mm to the left of the body block.
\usepackage{changepage}
\newenvironment{widebody}{%
  \begin{adjustwidth}{-16mm}{0mm}%
}{%
  \end{adjustwidth}%
}

% Full-bleed cover page support — image at exact paper size, title overlaid.
\usepackage{tikz}
\usetikzlibrary{positioning}

\begin{document}
\mainmatter
\silentchapter{Proposisjonar}{proposisjonar}

\prop{1}{}{Formverda er alt som er tilfelle.}
\silentchapter{2 Ikkje alle posisjonar er like sannsynlege}{ikkje alle posisjonar er like sannsynlege}
\prop{2}{}{Ikkje alle posisjonar er like sannsynlege.\footnote{Sannsynet er systemtilstandens vekt; ingen tenestekonfigurasjon flyt fritt.}}
\silentchapter{3 Seleksjonstrykka produserer eit landskap}{seleksjonstrykka produserer eit landskap}
\prop{3}{}{Seleksjonstrykka spenner ut eit landskap over formrommet.\footnote{Kompilatoren stettar syntaksen, køyretida allokerer minne, og nettet introduserer latens; summen av desse vilkåra spenner ut arkitekturens usynlege topografi.}}
\silentchapter{4 Landskapet er dynamisk}{landskapet er dynamisk}
\prop{4}{}{Landskapet er dynamisk.}
\silentchapter{5 Det finst agentar}{det finst agentar}
\prop{5}{}{Verda inneheld agentar som responderer på landskapet.\footnote{Ein ruter som allokerer bandbreidde og ein klient som navigerer eit grensesnitt er isomorfe traverseringar av det same topologiske rommet.}}
\silentchapter{6 Forma oppstår mellom agentane}{forma oppstår mellom agentane}
\prop{6}{}{Forma oppstår mellom agentane.\footnote{Arkitekturen er ikkje noko nodane eig, men tilstanden som trer fram i krysninga av anropa deira.}}
\silentchapter{7 Om det som ikkje lèt seg forme}{om det som ikkje lèt seg forme}
\prop{7}{}{Om det som manglar form, lyt ein teia.}

\addcontentsline{toc}{section}{A.1 Notasjon}
\prop{c}{}{ein klasse}
\prop{x}{}{ei form}
\prop{p}{}{eit seleksjonstrykk}
\prop{G}{}{ein grammatikk}
\prop{t}{}{eit tidspunkt}
\begin{overgang}π : Shape \msym{→} \msym{ℝ}ⁿ: morforomprojeksjonen\end{overgang}
\prop{H}{}{Shannon-entropi}
\prop{I}{}{gjensidig informasjon}
\addcontentsline{toc}{section}{A.3 Definisjonar}
\addcontentsline{toc}{section}{A.6 Empiriske testresultat}
Nærmaste-nabo-distansen i (H, W, D) etter z-skalering har ein variasjonskoeffisient (CV) på 5,4 (95 \% CI [3,7, 5,9]), noko som er om lag 15 gonger over Poisson-nullhypotesen på 0,36. Funna er robuste på tvers av 14 hold-out-subset (museum, periode, stil), inkludert NMK isolert (n = 63). n = 1664.

\clearpage
\addcontentsline{toc}{subsection}{A.6.1 Formrommet er ikkje uniformt busett (1.4)}
\par\addvspace{2pt}\noindent\centerline{%
  \includegraphics[width=\linewidth]{fig-A.6.1-uniformitet.pdf}}
\par\addvspace{2pt}
\anote{Vi måler gjensidig informasjon (sklearn k-NN-estimator) mellom kvar geometrisk dimensjon (Høgde, Breidde, Djupn, H/B) og to kandidatprediktorar: stilperiode og primærmateriale. Stilperiode forklarer 0.48–0.92 bits per dimensjon mot 0.23–0.42 bits for materiale (n = 1469). Den kulturelle merkelappen vinn over den fysiske eigenskapen med ein faktor 1.9\msym{×} til 2.3\msym{×} på alle fire dimensjonar samtidig. Funksjonalismen kunne ikkje predikert dette: materialet er den faktiske mekaniske avgrensninga; stilperioden ber ingen ergonomisk bodskap. Resultatet er, så vidt vi veit, den fyrste kvantitative testen av denne påstanden over eit fleirhundreårig korpus.}
\addcontentsline{toc}{subsection}{A.6.2 Stilperiode som samlevariabel (2.4, 2.62)}
\par\addvspace{2pt}\noindent\centerline{%
  \includegraphics[width=\linewidth]{fig-A.6.2-kanalisering.pdf}}
\par\addvspace{2pt}
\anote{Variasjonskoeffisienten over seks mesh-trekk (sphericity, complexity, fill-ratio, inertia-ratio, boks-volum, hylster-volum) strekkjer seg over to storleiksordnar (n = 2202). Sphericity er det mest kanaliserte trekket (CV = 0.074); hylster-volumet det friaste (CV = 9.52). Spreiinga er 128\msym{×}. Eit funksjonelt avgrensa formrom ville hatt jamn varians: kvart trekk skulle vere like fritt eller like bunde av tilpassingskrava. At nokre dimensjonar er låst og andre nær fri, er ein direkte målbar signatur av kanaliseringshierarkiet i tilpassingslandskapet.}
\addcontentsline{toc}{subsection}{A.6.3 Kanaliseringshierarki i mesh-rommet (3.3)}
\par\addvspace{2pt}\noindent\centerline{%
  \includegraphics[width=\linewidth]{fig-A.6.3-silhouette.pdf}}
\par\addvspace{2pt}
\anote{Silhuett-skoren for stilperiode i det 4-dimensjonale mesh-trekk-rommet (z-skalert) er −0.338, 95 \% CI [−0.346, −0.329] over 25 stilar med \msym{≥} 10 medlemar (n = 1971). Permutasjons-p < 0.001. Negativ silhouette tyder at gjennomsnittsstolen i ein stilperiode ligg nærmare nabostilane enn dei eigne medkandidatane. 21 av 25 stilar er negative; dei 4 positive er svært små samples (n = 12–16). Stilperiodar er kontinuerlege gradientar i morforommet, ikkje topologiske klynger. Vi har ikkje funne nokon tidlegare kvantifisering av denne påstanden over eit historisk møbelkorpus.}
\addcontentsline{toc}{subsection}{A.6.4 Stilar er gradientar, ikkje topologiske klynger (3.4)}
\par\addvspace{2pt}\noindent\centerline{%
  \includegraphics[width=\linewidth]{fig-A.6.4-ekspansjon.pdf}}
\par\addvspace{2pt}
\anote{Det kumulative konvekse hylsterveolumet i (H, B, D)-rommet, etter 1.–99. persentil-klipping per dimensjon, veks MONOTONT over alle tilgjengelege 50-årsperiodar (n = 1041 stolar med komplette dimensjonar). Frå 79 392 cm³ i fyrste periode til 589 796 cm³ ved slutten — ein 7\msym{×} ekspansjon utan ein einaste tilbakegang. Eit fast funksjonsenvelop ville krevja metning når korpuset veks. Den observerte monotone veksten støttar Kauffmans «det tilstøytande moglege» som aktivt ekspanderande rom og er, så vidt vi veit, fyrste gong dette er kvantifisert direkte for designartefakt.}
\addcontentsline{toc}{subsection}{A.6.5 Kumulativ ekspansjon av formrommet (4.4)}
\par\addvspace{2pt}\noindent\centerline{%
  \includegraphics[width=\linewidth]{fig-A.6.5-mahogni.pdf}}
\par\addvspace{2pt}
\anote{I utvalet av norskproduserte stolar frå perioden 1825–1849 brukar 16 av 16 mahogni. I den føregåande tilgjengelege 25-årsperioden (1750–1799) er fraksjonen 0 av 16. Lock-in skjer over éin generasjon. Mahogni og valnøtt er biomekanisk utbytbare for sittefunksjon: dei har samanliknbar tettleik, styrke og bearbeidbarheit. Ein binær kohort-overgang som dette kan ikkje forklarast som funksjonell optimisering — det er sti-avhengig seleksjonskollaps i sitt klaraste empiriske uttrykk.}
\addcontentsline{toc}{subsection}{A.6.6 Mahogni-kollaps 1825-1849 (4.5)}
\par\addvspace{2pt}\noindent\centerline{%
  \includegraphics[width=\linewidth]{fig-A.6.6-wasserstein.pdf}}
\par\addvspace{2pt}
\anote{Wasserstein-1 distansen mellom alle 10 suksessive 50-årsperiodar i korpuset gjev gjennomsnitt 15.8 cm for høgde, 9.2 cm for breidde og 5.8 cm for djupn. INGEN av dei 10 periodepara har distanse under terskelen 0.5 cm. Den statiske-landskap-nullhypotesen vert direkte falsifisert: kvar suksessiv periode har ei statistisk distinkt fordeling av former. Menneskeleg ergonomi har ikkje endra seg over 500 år; form endrar seg langt fortare enn funksjon.}
\addcontentsline{toc}{subsection}{A.6.7 Direkte falsifisering av postulat 4.1}
\par\addvspace{2pt}\noindent\centerline{%
  \includegraphics[width=\linewidth]{fig-A.6.7-trajektorie.pdf}}
\par\addvspace{2pt}
\anote{Sentroiden av kvar 50-årsperiode er ein punkt i (Breidde \msym{×} Høgde). Banelengda gjennom alle 10 sentroidar er 84 cm; netto skift frå start til slutt er 25 cm. Tortuositeten — banelengde delt på netto skift — er 3.45. Ein verdi over 1 tyder at banen vandrar fram og attende, ikkje monoton mot eit optimum. Tortuositet 3.45 er signaturen til eit dynamisk system med endrande seleksjonstrykk, ikkje av funksjonell konvergens (n = 1041).}
\addcontentsline{toc}{subsection}{A.6.8 Materiell kompleksitet per nasjon (5.3)}
\par\addvspace{2pt}\noindent\centerline{%
  \includegraphics[width=\linewidth]{fig-A.6.8-materialblanding.pdf}}
\par\addvspace{2pt}
\anote{For kvar stol med både materialliste og produksjonsland teiler vi talet på distinkte material i lista (n = 1582). Norske og danske stolar har medianverdi 3, britiske og italienske 2. Skilnaden held seg etter at vi har kontrollert for tid og stilperiode. Funksjonelt sett byggjer alle desse stolane den same tinga (eit sete for ein voksen menneskekropp); skilnaden er ein nasjonal-kulturell handverkssignatur. Vi kjenner ikkje til tidlegare arbeid som har kvantifisert dette over eit paneuropeisk korpus.}

\addcontentsline{toc}{subsection}{A.6.9 Materialstraumen 1500 til 2025 (4.5)}
\par\addvspace{2pt}\noindent\centerline{%
  \includegraphics[width=\linewidth]{fig-A.6.9-materialstraum.pdf}}
\par\addvspace{2pt}
\anote{Vi tokeniserer materiallista til kvar stol og normaliserer til éi av ti primærkategoriar. For kvar 25-årsperiode teljer vi kor mange stolar som inneheld kvart material. Resultatet er klare seleksjonsbølgjer som ingen før har lagt fram så detaljert: eik 1500–1700, nøttetre kring 1675, mahogni 1750–1850, modernismen sine material (stål, plast, kryssfiner, aluminium) etter 1900. Bølgjene er ikkje gradvise overgangar mellom funksjonelt overlegne alternativ — kvar bølgje er ein lokal kohort-event.}

\addcontentsline{toc}{subsection}{A.6.10 H/B-proporsjonen 1500 til 2024 (4.1)}
\par\addvspace{2pt}\noindent\centerline{%
  \includegraphics[width=\linewidth]{fig-A.6.10-proporsjon.pdf}}
\par\addvspace{2pt}
\anote{Den rullande 50-årsmedianen av høgde delt på breidde fell systematisk frå 1.88 i 1600 til 1.36 i 2000 (n = 1133). Endringa svarar til over eitt halvt standardavvik. Funksjonelt er proporsjonen 1.36 ikkje meir ergonomisk enn 1.88: begge er innanfor menneskeleg sitteområde. Det er ein form-drift, ikkje ein funksjonsdrift, som denne tidsserien direkte måler.}

\addcontentsline{toc}{subsection}{A.6.11 Materialnisjar i 3D-morforommet (5.3)}
\par\addvspace{2pt}\noindent\centerline{%
  \includegraphics[width=\linewidth]{fig-A.6.11-materialnisjar.pdf}}
\par\addvspace{2pt}
\anote{Kvar stol blir plotta som eitt punkt i (Breidde, Djupn, Høgde) og fargelagt etter primærmaterialet (n = 936). Sentroidane er klart åtskilde i den vertikale dimensjonen: trestolar (eik, nøttetre, mahogni) sit kring H \msym{≈} 85 cm, metall- og plaststolar kring H \msym{≈} 67 cm. Same fysiske oppgåve, klart ulike geometriske nisjar. Material er ikkje berre ein temporal merkelapp; det er ein geometrisk akse i morforommet.}

\addcontentsline{toc}{subsection}{A.6.12 Nyhetsraten og det tilstøytande moglege (6.5)}
\par\addvspace{2pt}\noindent\centerline{%
  \includegraphics[width=\linewidth]{fig-A.6.12-nyhetsrate.pdf}}
\par\addvspace{2pt}
\anote{Vi diskretiserer 3D-morforommet i 5 cm-voksler og teller kor mange voksler kvar 25-årsperiode opnar opp for fyrste gong. Nyhetsraten fell ikkje monotont mot metning slik ein lukka modell ville krevja; han hentar seg inn att i moderne tid (1925–1975, rate 0.5–0.6) når nye material kjem inn. Dette er ein direkte empirisk indikasjon på Kauffmans «det tilstøytande moglege» — ein mekanisme som ingen tidlegare har kunna måle på designartefakt.}

\addcontentsline{toc}{subsection}{A.6.13 Periodesentroidvandring i 3D (4.1)}
\par\addvspace{2pt}\noindent\centerline{%
  \includegraphics[width=\linewidth]{fig-A.6.13-3d-trajektorie.pdf}}
\par\addvspace{2pt}
\anote{Sentroiden av kvar 50-årsperiode er ein punkt i (Breidde, Høgde, Djupn). Banelengda i full 3D er 84 cm; netto skift frå start til slutt er 25 cm; tortuositeten er 3.45. Ein tortuositet over 1 tyder at sentroiden vandrar fram og attende — signaturen til eit dynamisk seleksjonslandskap, ikkje monoton funksjonell konvergens. Punktstorleiken kodar djupn, farge kodar tid (n = 1041).}

\addcontentsline{toc}{subsection}{A.6.14 Stilperiodefylogenese ved Ward-klynging (3.4)}
\par\addvspace{2pt}\noindent\centerline{%
  \includegraphics[width=\linewidth]{fig-A.6.14-fylogenese.pdf}}
\par\addvspace{2pt}
\anote{Vi reknar ut sentroiden av kvar stilperiode i 4D mesh-trekk-rommet og bygg ein hierarkisk klynge ved Ward-linkage. Resultatet er eit dendrogram som avdekker ein meiningsfull genealogi: rokokko sit nær barokk og renessanse, modernismen har sin eigen gren med funksjonalisme og Bauhaus, og 1800-tals stilane klynger saman. Funksjonen har ingen genealogi — berre forma kan arve frå ei tidlegare form. Dette er ein kvantitativ målbar form-arvs-struktur.}

\addcontentsline{toc}{subsection}{A.6.15 Rekurrensanalyse av periodesentroidar (4.1, 4.4)}
\par\addvspace{2pt}\noindent\centerline{%
  \includegraphics[width=\linewidth]{fig-A.6.15-rekurrens.pdf}}
\par\addvspace{2pt}
\anote{Symmetrisk avstandsmatrise over alle 19 25-årsperiodar i full (H, B, D). Mørke celler er like periodar; lyse er fjerne. Det modernistiske brotet etter 1900 viser seg som ei lys L-form i nedre høgre hjørne: ingen periode FØR 1900 liknar nokon periode ETTER 1900. Formhistoria gjentar seg ikkje på tvers av modernismen — eit nytt regime tek over. Vi har ikkje sett denne typen rekurrensanalyse brukt på designhistorie tidlegare.}

\addcontentsline{toc}{subsection}{A.6.16 Tilpassingslandskapet med stabile attraktorar (3.2)}
\par\addvspace{2pt}\noindent\centerline{%
  \includegraphics[width=\linewidth]{fig-A.6.16-fitnesslandskap.pdf}}
\par\addvspace{2pt}
\anote{KDE-tettleiken p̂(B, H) over alle stolar (n = 1681) har to klare lokale maksimum: éin ved (B \msym{≈} 50, H \msym{≈} 91) — den tradisjonelle høgrygga stolen — og éin ved (B \msym{≈} 50, H \msym{≈} 45) — den lågare modernistiske stolen. Dei tomme dalane mellom toppane er forbodne regionar i morforommet. Multimodaliteten er direkte støtte til proposisjon 3.2: tilpassingslandskapet har fleire stabile attraktorar, ikkje éin global optimum.}

\addcontentsline{toc}{subsection}{A.6.17 Endringsrate og episodisk diskontinuitet (4.3)}
\par\addvspace{2pt}\noindent\centerline{%
  \includegraphics[width=\linewidth]{fig-A.6.17-endringsrate.pdf}}
\par\addvspace{2pt}
\anote{Vi reknar ut Wasserstein-1 distansen mellom suksessive 25-årsperiodar i kvar av (Høgde, Breidde, Djupn) og aggregerer til ein samla rate R = \msym{√}(W\textsubscript{H}² + W\textsubscript{B}² + W\textsubscript{D}²) cm per 25-årssteg (n = 1105 stolar, 17 periodeparar etter krav om minst 10 stolar per bin). Sarle's bimodalitetskoeffisient er BC = 0,44, godt under terskelen 0,555. Streng bimodalitet er ikkje påvist: ratefordelinga er unimodal med ein lang hale. Men dei høgaste episodane er klart åtskilde frå mediananen 17,4 cm/25 år: 1650–1675 (34,0), 1825–1850 (28,7), 1675–1700 (28,3), 1800–1825 (26,2) og 1900–1925 (23,9). Dei stillaste periodane ligg under 13 cm. Proposisjon 4.3 brukar ordet «diskontinuerleg», ikkje «bimodal»: episodane stadfestar det. Det modernistiske brotet er tydeleg i ratane, men det 17. hundreårets stilskifte er endå tydelegare. «Diskontinuerleg» bør forståast som episodisk, ikkje som bimodalt fordelt.}

\addcontentsline{toc}{subsection}{A.6.18 Latente materialsignaturar (2.61)}
\par\addvspace{2pt}\noindent\centerline{%
  \includegraphics[width=\linewidth]{fig-A.6.18-materiallatens.pdf}}
\par\addvspace{2pt}
\anote{For kvar 50-årsperiode reknar vi ut sentroiden i (Høgde, Breidde, Djupn) for stolar som inneheld det aktuelle moderne materialet, og avstanden til sentroiden for stolar med utelukkande tre. Kvart panel viser materialets eige punkt-sky lagt over den same lyse bakgrunna av trestolar, slik at ein ser kor lenge det moderne materialet held seg innanfor det treaktige territoriet før det bryt ut. Stål viser signaturen klarast: avstand kring 10 cm i 1900 og 1950, så eit brått sprang til 38 cm i 2000 — ein latensperiode på minst femti år før divergensen. Plast og kryssfiner viser same forløp i mindre dramatisk form (begge under 13 cm i 1950, over 16 cm i 2000). Aluminium har berre eitt observerbart vindauge, 1950, der det ligg ekstremt nær trestolen (7,9 cm) — framleis i latent fase i korpuset. Alle fire moderne material er fyrst geometrisk uskiljelege frå tre, og dei som divergerer, gjer det brått. Materialet ber affordansen lenge før agentane realiserer han.}

\addcontentsline{toc}{subsection}{A.6.19 Det empiriske kraftfeltet i morforommet (3.1, 3.2, 3.3)}
\par\addvspace{2pt}\noindent\centerline{%
  \includegraphics[width=\linewidth]{fig-A.6.19-kraftfelt.pdf}}
\par\addvspace{2pt}
\anote{Tilpassingslandskapet er ikkje ein metafor — det er ei empirisk målbar topologi over morforommet. Vi reknar ut den gaussiske kjerne-tettleiken \msym{ρ̂}(B, H) over alle 1\,656 stolar med (Breidde, Høgde) og viser ho som ei fylt høgdekart med konturlinjer. Kraftfeltet \msym{−∇}V \msym{=} \msym{∇}log\,\msym{ρ̂} er overlagt som eit pilfelt på eit regulært rutenett: kvar piltipp peikar mot den næraste haugen i landskapet, slik agenten vil glide om ho fylgjer den lokale gradienten. Tre lokale maksimum vert oppdaga automatisk av eit maximum-filter og namngjevne av den hyppigaste stilperioden i kvar 9-cm-radius: Nyklassisisme (n = 361), Historisme (n = 314) og Modernisme/Midtjahrhundre (n = 170). Den turkise bana frå 1500 til 2025 er sentroiden av kvar 50-årsperiode i korpuset, lagt oppå landskapet som ei historisk vandring: tradisjonen sirklar lenge rundt den klassiske haugen, så glid han ned mot den modernistiske attraktoren etter 1900. Landskapet, gradientane, attraktorane og bana er alle utleidde frå dei same n stolane utan modell, fit eller imposed kategoriar. Proposisjonens sentrale diagram er for fyrste gong ein empirisk målt gjenstand.}

\addcontentsline{toc}{subsection}{A.6.20 Stolatlaset}
\par\addvspace{2pt}\noindent\centerline{%
  \includegraphics[width=\linewidth]{fig-A.6.20-stolatlaset.pdf}}
\par\addvspace{2pt}
\anote{Specimen-platet til appendikset: 1\,571 europeiske stolar i kronologisk orden, kvar teikna i sine eigne (Breidde, Høgde)-proporsjonar som ein eigen flis i mosaikken. Lesast spalte for spalte (40 spaltar \msym{×} 40 rader, sortert kolonnvis). Fyllfargen i kvar flis er primærmaterialet (tre, stål, kryssfiner, plast, aluminium, anna). Den indre kremfarga forma er stolens faktiske geometriske proporsjonar; ein tynn linje innanfor markerer setehøgda 45 cm; ein kort innvendig stripe peikar mot stolens avvik frå sentroiden av sin eigen 25-årskohort — kvart objekt ber sin individuelle signatur mot peer group. Frå avstand viser mosaikken materialepokane: tre dominerer 1315–1900 som ein nesten samanhengande brun mur; etter 1900 spring det ut ein eksplosjon av rust (stål), turkis (kryssfiner) og gull (plast). Frå nært hald les ein kvar flis som éin verkeleg gjenstand. Ingen har, så vidt vi veit, publisert eit slikt felt-i-felt-syn av eit europeisk møbelkorpus tidlegare.}

\backmatter
\silentchapter{Etterord}{etterord}
\begin{widebody}
Empirien er tufta på eit korpus av om lag 2 300 europeiske stolar (1280–2024). Datasettet tillet kvantitativ testing av formhistoriske hypotesar. Analysen av meshgeometri har synt at sju mesh-avleia trekk aukar den prediktive krafta for stilperiode med ein faktor på over fire samanlikna med katalogdimensjonar.

Overgangen frå prosa til proposisjonsform var ein logisk nødvendigheit. Desimalnummereringa angjev den logiske vekta, og superscript (d, a, t, o, i) indikerer logisk status. Til skilnad frå Wittgensteins original, er kvar proposisjon knytt til eit eksplisitt falsifiseringsvilkår; traktaten er konstruert for å kunne falsifiserast. Dei empiriske testane stadfestar hovudproposisjonane: formrommet er ikkje-uniformt busett; stilperiode er ein effektiv proxy for aggregerte seleksjonstrykk; og landskapet er i uopphøyrleg endring. Funna er robuste på tvers av ulike geometriske representasjonar.

Observert latens i stålillustrasjonen viser at formhistoria er langsamare enn materialhistoria; eit nytt substrat arvar formspråket til det substratet det erstattar, heilt til eigne affordansar pressar det ut i nye regionar.

Traktaten skildrar krefter og posisjonar, ikkje verdiar. Ho kan forklare avstanden i formrommet, men ho kan ikkje felle estetiske domar. Estetikk er transcendent til formsystemet, på same vis som etikk er transcendent til logikken (Wittgenstein). Agenten navigerer; modellen registrerer berre den resulterande posisjonen. Proposisjonane skal erkjennast som ein stige. Dei skal ikkje erstatte handverket, men gjere den intuitive navigasjonen eksplisitt ved å vise agenten kva han allereie gjer.

At teksten kan falsifiserast, er garantien for hennar gyldigheit. Om eitt einaste seleksjonstrykk forklarer all formvariasjon, kollapsar den logiske kjeden. Traktaten er stigen som skal kastast.

Takk til Jan Petter Neverdahl og Hermann Stange for kommentarar under arbeidet.

\end{widebody}


\silentchapter{Notatar}{notatar}
\begin{widebody}
\theendnotes
\end{widebody}

\silentchapter{Referansar}{referansar}
\begin{widebody}
\refentry{Sullivan, L. H. (1896). The tall office building artistically considered. Lippincott's Magazine, 57, 403–409.}
\refentry{Thompson, D'A. W. (1917). On growth and form. Cambridge University Press.}
\refentry{Wright, S. (1932). The roles of mutation, inbreeding, crossbreeding, and selection in evolution. Proc. Sixth Int. Congress of Genetics, 1, 356–366.}
\refentry{Rosenblueth, A., Wiener, N. \& Bigelow, J. (1943). Behavior, purpose and teleology. Philosophy of Science, 10(1), 18–24.}
\refentry{Shannon, C. E. (1948). A mathematical theory of communication. Bell System Technical Journal, 27, 379–423.}
\refentry{Wiener, N. (1948). Cybernetics. MIT Press.}
\refentry{Turing, A. M. (1950). Computing machinery and intelligence. Mind, 59(236), 433–460.}
\refentry{Waddington, C. H. (1957). The strategy of the genes. George Allen \& Unwin.}
\refentry{Kubler, G. (1962). The shape of time. Yale University Press.}
\refentry{Kuhn, T. S. (1962). The structure of scientific revolutions. University of Chicago Press.}
\refentry{Raup, D. M. (1966). Geometric analysis of shell coiling. Journal of Paleontology, 40(5), 1178–1190.}
\refentry{Eldredge, N. \& Gould, S. J. (1972). Punctuated equilibria. I T. J. M. Schopf (Red.), Models in paleobiology, 82–115.}
\refentry{Stiny, G. \& Gips, J. (1972). Shape grammars and the generative specification of painting and sculpture. Information Processing 71, 1460–1465.}
\refentry{Gibson, J. J. (1979). The ecological approach to visual perception. Houghton Mifflin.}
\refentry{Stiny, G. (1980). Introduction to shape and shape grammars. Environment and Planning B, 7(3), 343–351.}
\refentry{Stiny, G. (1991). The algebras of design. Research in Engineering Design, 2, 171–181.}
\refentry{Kauffman, S. A. (1993). The origins of order. Oxford University Press.}
\refentry{Arthur, W. B. (1994). Increasing returns and path dependence in the economy. University of Michigan Press.}
\refentry{Michl, J. (1995). Form follows WHAT? The modernist notion of function as a carte blanche.}
\refentry{Odling-Smee, F. J., Laland, K. N. \& Feldman, M. W. (2003). Niche construction. Princeton University Press.}
\refentry{Hatchuel, A. \& Weil, B. (2003). A new approach of innovative design: An introduction to C-K theory. Proc. ICED 2003.}
\refentry{Hatchuel, A. \& Weil, B. (2009). C-K design theory: An advanced formulation. Research in Engineering Design, 19(4), 181–192.}
\refentry{Stiny, G. (2006). Shape: Talking about seeing and doing. MIT Press.}
\refentry{Mitteroecker, P. \& Huttegger, S. M. (2009). The concept of morphospaces in evolutionary and developmental biology. Biological Theory, 4(1), 54–67.}
\refentry{Fields, C. \& Levin, M. (2022). Competency in navigating arbitrary spaces. Entropy, 24(6), 819.}
\refentry{Levin, M. (2022). Technological approach to mind everywhere. Frontiers in Systems Neuroscience, 16, 768201.}
\refentry{Levin, M. (2025). Living things are not machines (also, they totally are). Noema Magazine.}
\end{widebody}

\end{document}
